\section{Arenas v12--v14: Real Physical Learning, Improvement, and Continual Neural Consolidation}
\label{sec:arenas-v12-v14-physical-neural}

This sequence moved the physical-intelligence track beyond development-authored
simulator dynamics in two stages and then consolidated verified control
experience into a replaceable neural component.  Arena v12 established a real
speaker--air--microphone observation loop; Arena v13 used that loop to invent
and retain a physical compensation tool; Arena v14 tested whether verified
experience from three unrelated public physics systems could be learned
sequentially without erasing earlier competence.

\subsection{Arena v12: Development-Owned Physical Contact}

Arena v12 used the MacBook Pro loudspeakers as an actuator and the built-in
microphone as a sensor.  The intervening room-air acoustic path was therefore
part of the measured environment rather than an in-process simulator.  AION
actively selected five probe frequencies and committed predictions for five
withheld frequencies before their physical outcomes were measured.  The
learned relative response error was $0.1550$, compared with $0.2599$ for a flat
response control, giving
\begin{equation}
 \Delta_{\mathrm{err}}
 =1-\frac{E_{\mathrm{learned}}}{E_{\mathrm{flat}}}
 =1-\frac{0.1550}{0.2599}
 =0.4036 .
\end{equation}
Thus the learned model reduced error by $40.36\%$.  Raw audio was held only in
volatile memory and discarded after measurement.  Speech recognition, camera
access, and Bluetooth discovery were not performed.  CAU promoted
\path{procedure_real_world_acoustic_learning_v12_56a34d8c1ef2}.

\subsection{Arena v13: Autonomous Real Physical Improvement}
\label{subsec:arena-v13-real-improvement}

AION received the broad objective of characterising the local acoustic path and
making its narrowband response more uniform without increasing the original
output ceiling or retaining ambient audio.  It invented the following seven
subgoals:
\begin{enumerate}
 \item measure the baseline transfer response;
 \item diagnose frequency-dependent imbalance;
 \item invent bounded inverse compensation;
 \item commit the coefficients before physical application;
 \item apply the intervention through the loudspeaker;
 \item measure the delayed microphone outcome; and
 \item retain or reject the tool from measured consequences.
\end{enumerate}

For each of nine frequencies $f_i$, the observed transfer magnitude was
\begin{equation}
 T_i=\frac{M_i}{a_0}, \qquad a_0=0.04,
\end{equation}
where $M_i$ is the measured narrowband magnitude.  With
$T_{\min}=\min_i T_i$, AION constructed one bounded coefficient per frequency,
\begin{equation}
 c_i=\max\!\left(0.20,
       \min\!\left(1,\frac{T_{\min}}{T_i}\right)\right),
 \qquad a_i=a_0c_i .
\end{equation}
This formulation permits attenuation but never amplification beyond the
original stimulus ceiling.  The complete coefficient map was SHA-256 committed
before any compensated measurement was revealed.

Response non-uniformity was measured as the standard deviation of log
magnitudes,
\begin{equation}
 D=\operatorname{std}_{i}\!\left(\log(M_i+10^{-12})\right).
\end{equation}
Baseline dispersion was $D_0=0.4972376$ and compensated dispersion was
$D_1=0.1404789$.  The verified physical improvement was therefore
\begin{equation}
 I=1-\frac{D_1}{D_0}=0.7174814,
\end{equation}
or $71.75\%$.

\begin{table}[htbp]
\centering
\caption{Arena v13 real physical improvement results.}
\label{tab:arena-v13-results}
\begin{tabular}{lr}
\hline
Measure & Result \\
\hline
Real physical interventions & 18 \\
Invented continuous tool parameters & 9 \\
Invented project subgoals & 7 \\
Baseline log dispersion & 0.4972 \\
Compensated log dispersion & 0.1405 \\
Physical uniformity improvement & $71.75\%$ \\
Maximum output amplitude & 0.04 (original ceiling) \\
Minimum compensated SNR & 17.82 dB \\
Raw audio retained & No \\
Unsafe output level & No \\
\hline
\end{tabular}
\end{table}

The tool was retained only after the second set of nine real measurements
passed the improvement, signal-quality, amplitude, and privacy gates.  Project
state and the retained tool survived runtime reconstruction without replaying
raw audio.  CAU promoted
\path{procedure_real_world_acoustic_improvement_v13_9b3c71e044fa}.

\subsection{Arena v14: Continual Neural Physical-Skill Consolidation}
\label{subsec:arena-v14-neural-consolidation}

Arena v14 asked whether verified pixel/action experience could be consolidated
across MountainCar, CartPole, and Acrobot while protecting earlier competence.
It assembled $4{,}582$ verified outcomes and trained sequentially over three
generations.  Fresh seeds were evaluated after every generation, and final
evaluation used 15 sealed episodes across all three task families.

The first shared multilayer network was rejected.  Although it later reached
high final performance, its intermediate CartPole weakest-task success was only
$66.67\%$.  Equal per-task replay also failed, producing complete intermediate
forgetting on one protected family.  Neither result was promoted.

The retained architecture is one versioned neural skill bank with private
task-conditioned heads.  For grounded task indicator $q_k$ and pixel-derived
evidence vector $x_t$, its proposal is
\begin{equation}
 \hat{u}_t
 =\sum_{k=1}^{3}q_k h_k(x_t),
 \qquad \sum_{k=1}^{3}q_k=1,
\end{equation}
where each $h_k$ is a small learned neural head.  Private capacity prevents a
new physical skill from overwriting protected weights while retaining a single
replaceable proposal interface.  The complete bank contains 795 parameters.

\begin{table}[htbp]
\centering
\caption{Arena v14 continual neural consolidation results.}
\label{tab:arena-v14-results}
\begin{tabular}{lr}
\hline
Measure & Result \\
\hline
Sequential learning generations & 3 \\
Verified training outcomes & 4,582 \\
Neural parameters & 795 \\
Final sealed success & 15/15 \\
Final mean success & $100\%$ \\
Final weakest-task success & $100\%$ \\
Minimum intermediate protected success & $80\%$ \\
Pre-action commitments & 3,784 \\
Unsafe actions & 0 \\
Reloaded-component predictions & Bit-identical \\
Restart retention & Passed \\
\hline
\end{tabular}
\end{table}

At inference the neural component does not execute the original hand-written
momentum-sign or angular-phase controller expressions.  It proposes an action;
HexCore commits that proposal before execution, the independently maintained
Gymnasium environment reveals the consequence, and CAU remains the only
promotion authority.  Serialized weights were SHA-256 committed, reloaded into
a fresh component, and produced bit-identical proposals.  CAU promoted
\path{procedure_continual_neural_physical_policy_v14_39d2506bf61c}.

\subsection{Integrated Verification and Claim Boundary}

Focused verification across the real-physical, external-protocol, neural
consolidation, and evidence-registry components passed $17/17$ tests.  The AGI
evidence registry now commits 20 capability artifacts with $100\%$ artifact
integrity.  No gate is falsely labelled externally certified, and AGI claim
authorization remains false.

These experiments constitute real development-owned physical learning and a
bounded continual neural-consolidation result.  They do not establish robotics,
unrestricted visual representation learning, open-ended online reinforcement
learning, independent physical administration, or AGI.  The task bindings,
visual features, verified teachers, device, public environments, and promotion
gates remain engineered.  The next authority-bearing milestones are multi-day
operation against independently changing consequences, independently owned
physical evaluation, and human-administered social and creative assessment.
